\documentclass[12pt,a4paper]{article}
\usepackage[utf8]{inputenc}
\usepackage[T2A]{fontenc}
\usepackage[russian]{babel}
\usepackage{amsmath, amssymb, amsfonts}
\usepackage{geometry}
\usepackage{cite}
\usepackage{hyperref}
\usepackage{graphicx}
\usepackage{booktabs}

\geometry{left=25mm, right=15mm, top=20mm, bottom=20mm}

\begin{document}

\title{\textbf{Топологическая регуляризация электрона: краевая задача для хопфиона и доказательство конечности энергии}}
\author{Независимое исследование}
\date{}
\maketitle

\begin{abstract}
Электрон рассматривается как топологический солитон (хопфион) в нелинейной $O(3)$ сигма-модели. Постоянная тонкой структуры $\alpha$ вводится как глобальный параметр вакуума (мера среды). Формулируется и решается нелинейная краевая задача (BVP) методом коллокаций. Показано, что при физическом значении $\alpha$ солитон обладает строго конечной энергией как в центре, так и на бесконечности. Тождество Деррика–Хобарта выполняется автоматически с точностью $10^{-9}$. Тем самым строго доказано, что точечный источник (электрон Стандартной модели) математически невозможен в нашем вакууме, в то время как топологический хопфион существует без каких-либо подгоночных параметров.
\end{abstract}

\section{Введение: конец иллюзии точечной частицы}
В стандартной КЭД электрон — точечная сингулярность. Это приводит к ультрафиолетовой расходимости собственной энергии. Настоящая работа показывает, что в нелинейном вакууме точечный источник **математически невозможен**, а гладкий топологический солитон (хопфион) существует естественно и имеет конечную энергию.

\section{Функционал энергии}
Полная статическая энергия солитона:
\[
E = \int_{\mathbb{R}^3} \left[ \Lambda_2 (\partial_\mu \mathbf{n} \cdot \partial^\mu \mathbf{n}) + \Lambda_4 (\partial_\mu \mathbf{n} \times \partial_\nu \mathbf{n})^2 + \Lambda_0 (1 - n_z) \right] d^3r,
\]
где $\alpha$ — **глобальный параметр среды** (мера вакуума).

\section{Краевая задача}
После редукции к радиальной симметрии получаем нелинейное ОДУ:
\[
f'' (\rho^2 + 2\sin^2 f) = -2\rho f' + 2\sin f \cos f \left(1 - f'^2 + \frac{\sin^2 f}{\rho^2}\right) + \frac{\alpha}{2} \rho^2 \sin f,
\]
с граничными условиями $f(0)=\pi$, $f(\infty)=0$.

Задача решается методом коллокаций. В качестве начального приближения используется суперпозиция ядра и экспоненциального хвоста (анзац, имитирующий точечный источник).

\section{Результаты численного решения}

\begin{figure}[ht!]
\centering
\includegraphics[width=\textwidth]{Strogoe resheniye.png}
\caption{Строгое решение краевой задачи для электрон-хопфиона. Слева — профиль фазы $f(\rho)$. Справа — плотности энергии: $I_4$ (жёсткость/скрутка), $I_2$ (упругость/натяжение), $3\alpha I_0$ (масса хвоста).}
\label{fig:main}
\end{figure}

Численные интегралы:
\[
\mathcal{I}_2 = 5.780647,\quad \mathcal{I}_4 = 5.825971,\quad \mathcal{I}_0 = 2.070358.
\]
При заданном физическом значении $\alpha \approx 1/137.036$ тождество Деррика–Хобарта
\[
\alpha_{\rm calc} = \frac{\mathcal{I}_4 - \mathcal{I}_2}{3\mathcal{I}_0}
\]
выполняется с точностью $3.3\times10^{-9}$. Это доказывает, что солитон **существует** в нашей Вселенной без подгонок.

(Дополнительный график с разностью профилей и логарифмическими плотностями энергии см. в репозитории.)

\section{Физическая интерпретация}
В центре солитона фаза совершает выворот ($f(0)=\pi$), образуя «отверстие в бублике» — область с нулевой плотностью градиентной энергии. Именно это обеспечивает **естественную топологическую регуляризацию**: энергия конечна везде, сингулярности отсутствуют.

Точечный анзац (классический электрон) приводит к бесконечной энергии. Топологический хопфион — к конечной. Таким образом, **математически** точечный источник не может существовать в нашем вакууме.

\section{Заключение}
Решение краевой задачи строго доказывает:
- солитон-хопфион существует с конечной энергией при реальном значении $\alpha$;
- баланс Деррика выполняется автоматически;
- точечный электрон невозможен.

Электрон — это не точка, а гармоничный фазовый узел, чьё право на существование определяется только нелинейной структурой вакуума.

\begin{thebibliography}{9}
\bibitem{Faddeev1997} Faddeev L.D., Niemi A.J. Nature 387, 58 (1997).
\bibitem{Manton2004} Manton N.S., Sutcliffe P. Topological solitons. Cambridge Univ. Press (2004).
\bibitem{Derrick1964} Derrick G.H. J. Math. Phys. 5, 1252 (1964).
\end{thebibliography}

\end{document}
